\label{sec:discussion}

{\revstart \emph{Security.} Inside a marked region, an intended origin and an eligible unlisted origin receive the same \openscope{} state. That state records two facts: the route falls below a holder-signed marker, and its length stays within the policy bound. It does not record that the origin is authorized, and an exact ROA still takes precedence. Section~\ref{sec:security-scope} states the adversary model, and Section~\ref{sec:eval-security} measures the reachability the state creates. The operator-facing decision is therefore narrow: which regions to mark, and which local action to apply. Because the state is reported explicitly, a router can count or log \openscope{} admissions separately from authorized routes and revisit the region once the traffic pattern settles.}

{\revstart \emph{What the conflict rates express.} At 80\% policy coverage, the relative reduction approaches the coverage itself, because a marked aggregate contributes no descendant conflict unless its origin is ineligible. The 79.1\% figure is therefore an internal-consistency check that states the size of the uncovered remainder, 10.5\% and 5.2\%, rather than an independently estimated effect. The replay uses 589 route--origin pairs under 97 aggregate prefixes from a single RRC00 snapshot. It characterizes that observed set, and it does not estimate Internet-wide rates. Reporting the absolute residual rates alongside the relative change keeps that distinction visible.}

{\revstart \emph{What the routing model supports.} The model fixes one 77,429-AS graph, one 589-AS upgrade set, 200 sampled units, and five upgrade orders. The intended-only and competition panels isolate benefit from exposure, but they remain separate conditional comparisons rather than a same-scenario frontier. The IPv6 sensitivity likewise leaves the median intended-delivery change at 0.0. Traffic, operator intent, another topology, and another ROV set all remain outside the evidence. The two panels answer separate questions: whether the intended route is reachable, and what else becomes reachable beside it.}

{\revstart \emph{Deployment path.} Holders place nonoverlapping markers before a matching ROA exists, operators map \openscope{} to a chosen local action, and exact ROAs authorize the origin once it is known. On the observed route set, prioritizing aggregates by the eligible route--origin pairs they cover reduces conflicts faster than ASRank, AS degree, or random ordering. A holder narrows a marked region by replacing one broad marker with disjoint markers on the intended subtrees. Once the matching ROA is in force, revoking the marker closes the region without touching the aggregate entry.}
